\chapter{Introduction - How to use this template}

\section{Choose your stile}

Since no style reference has been provided, we are quite free: you have first of all to choose the style and change some lines in the file \url{thesis.tex}. When two options are available, one of the two is commented starting the line with \%. To change it, comment what you don't need and de-comment the other option:

\subsection{Single- or Double-sided}
I like better the Single-sided, but if you want to go double-sided remove at line 14
\begin{verbatim}
\documentclass[12pt,a4paper,openright,oneside]{book}	%%One sided
\end{verbatim}
and de-comment line 15
\begin{verbatim}
\documentclass[12pt,a4paper,openright,twoside]{book}	%%Double sided
\end{verbatim}

\subsection{Headers and footers}
Lines 22-42 are going to use 


\subsection{PDF settings}
Insert the information about you and the thesis in lines 56-61. They will be included in the PDF information
\begin{verbatim}
\hypersetup{
    pdfauthor={AuthorName},
    pdftitle={shortTitle},
    pdfsubject={subject},
    pdfkeywords={keyword1, keyword2}
}
\end{verbatim}

\subsection{Language}
The main language is the last, so for a thesis in english line 69 is
\begin{verbatim}
\usepackage[italian,english]{babel} 
\end{verbatim}
Swap ita with eng fi you are writing in italian as main language




\section[Sectioning]{Chapters, sections and subsections}

A thesis has three parts: head, body and tail. Three correspondent folders are containing the source files for these parts.

Any of this source files is imported in the \url{thesis.tex} file, which is the only one to get compiled. This is realised using the command
\begin{verbatim}
\input{subfolder/filename.tex}
\end{verbatim}
To exclude a section from the compiling process, comment or remove the correspondent line. If no subfolder is specified, \LaTeX\ looks for it in the same folder of \url{thesis.tex}

E.g. the introduction source file is in the body subfolder, and is included using the command \begin{verbatim} \chapter{Introduzione}
\label{chap:first-chapter-intro}
Il campo dell'Anomaly Detection sta ricevendo molto interesse nell'ultimo periodo. L'avanzamento tecnologico ha portato alla nascita dei Big Data e degli approcci Data Driven consentendo a tecniche di Machine Learning di spopolare in tantissimi ambiti, tra cui quello dell'industria 4.0.
Attraverso l'innovazione degli impianti di produzione e all'utilizzo di dispositivi IoT, gli operatori possono monitorare il comportamento di un canale di produzione e intervenire nel caso di anomalie oppure pianificare la produzione grazie all'analisi di dati statistici o report.
Solamente questo ha permesso un'evoluzione sostanziale ma non ci si ferma qui, l'enorme quantità di dati raccolti rende possibile l'applicazione di tecniche di Machine Learning come supporto alle scelte decisionali in diversi ambiti: Anomaly Detection, Predictive Maintenance, previsioni e molto altro.

L’obiettivo di questa tesi é quello di contribuire al progetto Beat 4.0 portato avanti da SKF e ALTEN ITALIA. SKF é una azienda multinazionale specializzata nella produzione di cuscinetti a sfera, ALTEN é un'azienda leader nel mercato europeo della consulenza informatica e che sta offrendo supporto ad SKF. Questa collaborazione ha lo scopo di portare innovazione digitale all'impianto SKF di Cassino in provincia di Frosinone, attraverso l'applicazione di tecniche di Machine Learning.
Diverse collaborazioni sono state concluse, mentre altre ancora in corso, tra ALTEN e l'Università degli Studi di Torino per portare avanti il progetto Beat 4.0. Ognuna di queste si focalizza in un dominio specifico del Condition Monitoring: previsione della qualità, studio di correlazione tra i sensori e la qualità finale, Predictive Maintenance ed infine Anomaly Detection su cui si basa il lavoro di questa tesi.
Il contributo dello studente si concentra su tutto il ciclo di vita di produzione di tecniche di Machine Learning: studio qualitativo e processamento dei dati, analisi e applicazione dei metodi di Anomaly Detection ed infine distribuzione degli algoritmi su Cloud Azure e sviluppo front-end attraverso PowerBI.
Il focus della tesi si é concentrato prevalentemente sullo studio e sull'applicazione di tecniche di Model Selection non supervisionato: la mancanza di etichette all'interno del dataset fornitoci da SKF non permette la valutazione di un modello di Anomaly Detection, costringendo quindi a ricercare metodi che non ne richiedessero l'uso. Solitamente il Model Selection si effettua andando a partizionare i dati in modo diverso per la fase di train-test, oppure facendo ottimizzazione degli iperparametri, ma in ogni caso la presenza di etichette é necessaria. Per questo motivo vengono introdotte due tecniche di Model Selection non supervisionate, chiamate Meta Learning e Metriche Surrogate. La prima richiede di allenare dei rilevatori di anomalie su dataset etichettati per poi trasferire le informazioni apprese ad un dataset non etichettato che sia sufficientemente simile al primo, quindi questo approccio é stato scartato. Il secondo si basa sulla creazione di metriche non supervisionate, chiamate surrogate, che siano sufficientemente correlate alle più classiche metriche supervisionate come F1-Score. Verranno quindi proposte tre metriche surrogate chiamate \textit{Model Centrality}, \textit{Clustering Coefficient} e \textit{Performance on Injected Synthetic Anomalies} insieme a diversi metodi di Rank Aggregation: Borda, Robust Borda, AVG Score e Kemedy-Young utilizzati per combinare insieme le 3 metriche non supervisionate.

Il proseguo della tesi si divide quindi in questo modo: i \textit{Capitoli \ref{chap:intro}}, \textit{\ref{chap:methods}} e \textit{\ref{chap:modelselection}} presentano un'introduzione dell'Anomaly Detection, metodi di Anomaly Detection e Thresholding, e Model Selection rispettivamente. Nel \textit{Capitolo \ref{chap:skf}} vengono trattati i dati SKF e le loro proprietà e caratteristiche. Il \textit{Capitolo \ref{chap:impl}} tratta tutta la fase di sperimentazione, con una prima parte dedicata all'implementazione delle metriche surrogate e di Rank Aggregation, una seconda parte contenente i risultati degli esperimenti su dataset di benchmark provenienti da ODDS e SMD ed infine una terza parte in cui viene testato l'algoritmo di Model Selection su SKF. Il \textit{Capitolo \ref{chap:deploy}} conclude il ciclo di vita del software andando a trattare la fase di pubblicazione su Cloud Azure.
Conclusione e sviluppi futuri sono invece trattati nel \textit{Capitolo \ref{chap:conclusion}}. \end{verbatim}

\subsection{Front page}

Two front pages are provided: 
\begin{enumerate}
\item \url{frontPage.tex} for thesis with a single Relatore
\item \url{frontPage-cr.tex} for thesis with Relatore and Correlatore
\end{enumerate}
choose what is needed adapting lines 94-95.

\subsection{Dedication}
Dedication is the initial inspirational quote. Can be edited in the file \url{dedication.tex} or removed.

\subsection{Chapter names}
Chapter and section names are in this form:
\begin{verbatim}
\chapter[Name in the index]{Name on the title}
\end{verbatim}
To exclude them from the numbering, use the *; If no name for the index is specified, will be assumed the same. For example:
\begin{verbatim}
\chapter*{Name on the title}
\end{verbatim}
won't be numbered and will have the same name in the title and in the index.

\section{Bulleting}

Bullet list example
\begin{itemize}
\item first point
\item second point
\item third point \\
\end{itemize}

\noindent
Enumeration example
\begin{enumerate}
\item first point
\item second point
\item third point\\
\end{enumerate}

\noindent
Description example
\begin{description}
\item[first descr] first point
\item[second descr]  second point
\item[third descr]  third point\\
\end{description}

\noindent
\dots but you can also build nested lists
\begin{itemize}
\item first point
	\begin{itemize} %tab is only for the readability, but not necessary
	\item first point
	\item second point
	\end{itemize}
\item second point
\item third point \\
\end{itemize}


\section{Floating objects}

Floating objects are tables, figures and so on. 

\subsection{Position, captions and cross-referencing}

To force the position of an object use the options, for example:
\begin{verbatim}
\begin{figure}[htb]
\end{verbatim}
where [htb] means that the priority of positioning is \url{h} = here, \url{t} = top, \url{b} = bottom. Also \url{p} can be used and means that the object is placed in a page on itself.

Captions are managed automatically using the command
\begin{verbatim}
\caption{The title of my object}
\end{verbatim}
to cross reference to the object use the 
\begin{verbatim} 
\label{labelName} 
\end{verbatim} 
to assign a name to the object and  
\begin{verbatim} 
\ref{labelName}
\end{verbatim} 
to refer to that.



\subsection{Figures}

Insert a figure using the code
\begin{verbatim} 
\begin{figure}[htb]
\centering
\includegraphics[scale=0.15]{pictures/logo.png}
\caption{the logo of UniTo}
\label{myFigure}
\end{figure}
\end{verbatim} 

then you can cross-reference to it. In this case the discussion is about Fig. \ref{myFigure}.

Where the number of the figure has been gotten using 
\begin{verbatim} 
\ref{myFigure}
\end{verbatim} 


\begin{figure}[htb]
\centering

\includegraphics[scale=0.15]{pictures/logo.png}
\caption{the logo of UniTo}

\end{figure}

\subsection{Subfigures}

\begin{figure}[htb]
\centering
\subfigure[title of first subfigure]
	{\includegraphics[scale=0.15]{pictures/logo.png}}
\subfigure[title of second subfigure]
	{\includegraphics[scale=0.3]{pictures/logo.png}}	
\caption{Two logos of UniTo of different sizes}

\end{figure}

\subsection{Tables}
